\section{2026 年 11 月 1 日：study\_18 恢复与片段容错}
首轮 study\_18 的第 52 号片段返回人物/实体片段，而不是完整情景记忆，
被校验器拒绝，导致后续索引和 14 题问答未执行。
本次在首轮数据库的独立克隆上定向恢复，不覆盖首轮数据库或成绩。

恢复仅新增 1 条记忆，复用的 83 条正文及结构化内容逐项不变，
最终完成 84/84 条记忆和索引；14/14 题有效回答，无超时和评委错误。
恢复视频参考均分为 \textbf{0.3929}。
保持原四视频的答案与评分不变，仅补齐该视频后，
65 题组合均分为 \textbf{0.5538}，有效答案 59/65。
这是覆盖恢复，不是对所有题重跑后择优，也不是模型或记忆架构的独立增益；
首轮 0.4692 的记录仍完整保留。共同评委下的 M3-Agent 对照属于另一次评分。

oGMemory 的改动已提交到 \code{m3_dev_1031}：
\code{99b22866} 保存模型推理参数和有界并发，
\code{40ac1dcd} 增加无效输出重试及显式部分覆盖，
\code{c992d45b} 保证连接指纹不包含凭据及其哈希。
域无效输出可额外重试一次；若仍有少量失败，
只有至少 95\% 覆盖、回执与图一致、且没有旧失败节点可能混入时，
才允许索引与问答继续。缺失片段进入提示和每题回执；
全量完成的续跑判定不会误接受部分证据。
鉴权、存储或低覆盖继续失败关闭。

相关扩展测试为 1063 passed、1 deselected；排除的一项音频重建重复节点测试
已在修改前源码复现，不属于本次引入。另有独立 SQL 故障注入：
强制片段 52 连续失败，剩余 83 条仍实际建索引，测试 QA 继续并披露缺失，
最终退出码保持 2。该测试使用 stub 回答，不计入模型成绩。
